\section{Relation to earlier constructions}
\label{sec:finite-prior-comparison}
The determinant $\det[1,v,\ldots,v^{d-1}]$ is a classical local
index form. Arpin, Bozlee, Herr and Smith construct schemes of
monogenerators and polygenerators for finite locally free algebras
\cite{ABHS}. Their Proposition 3.6 and Definition 3.12 identify the
index-form open and its closed complement; Proposition 3.14 gives
the polygenerator open by minors. Thus neither the generator functor
nor its index determinant is claimed as new here. Our use concerns
the coherent cokernel of a global section-multiplication map and,
for generating hyperplanes, its exact closed incidence with a
rank-two quotient. Lemma~\ref{lem:hyperplane-subalgebra} and
Theorem~\ref{thm:universal-hyperplane} prove that closed-scheme
identification on nonreduced bases, rather than deriving it from
agreement of generator opens.

The algebra $G_r$ is the classical two-row Grassmannian cohomology
ring; \cite[Theorem 2.7]{Grinberg} provides the rectangular Schur
basis in a more general quotient. Its role here is an exact
transverse algebra of multiplication failure, determined by a
rank-two quotient of $\C[z]/(z^r)$. The explicit residual equations
\eqref{eq:hyperplane-residual-generators} verify this identification
without a tangent-cone approximation. The diagonal-ideal and
all-powers identities in the three-plane result are precisely
\cite[Corollary 3.8.3]{Haiman}; the argument here supplies the
finite total-degree presentation and conductor transport, not a
new proof of that theorem.

Normal generation in sufficiently positive degree on a curve is
classical \cite{GL86}. The curve conductor theorem uses it only for
$L(-D)$, and then proves the precise quotient comparison by filling
the double-vanishing space and its first normal layer. On $\Pj^1$
the evaluation-kernel argument gives the sharper rank-dependent
bound and its boundary example.

For pencils, confluent Vandermonde factorization and the discriminant
identity are classical algebra. The moving normalization is the
explicit scalar-pair incidence of that index-form divisor. Its
finite map and differential test are proved in
Section~\ref{sec:moving-contact-normalization}; the normalization
is not offered as a construction independent of index-form geometry.
The contact-hyperplane incidence is different: it is the failure
scheme itself, not just its normalization, and its total space is
smooth even at collisions.


\subsection*{Subalgebra varieties, quotient flags, and the additional comparisons}
Iovanov and Sistko \cite[Theorem 0.1]{IS} classify maximal subalgebras
of finite-dimensional algebras over fields under their Schur hypothesis.
In the commutative algebraically closed case, the alternatives include
identifying two residue factors and imposing a first-order condition
at one factor. These closed-point mechanisms are classical and are
not claimed here as new.
Sistko \cite[Definition 3.1, Theorems 1.1 and 4.1]{Sistko} constructs
subalgebra varieties and computes the maximal-subalgebra vanishing
ideal for basic algebras. The distinction needed here is between
such a variety and the full relative quadratic-equation functor:
for $\C[z]/(z^r)$ the latter already has the transverse algebra
$G_r$, of length $\binom r2$. Agreement on the reduced variety does
not identify this thickening.

The Hilbert functor used in Section~\ref{sec:quotient-hilbert} is the
standard functor of locally free quotient algebras; its representability
is not a contribution of this article. Likewise, nested effective
divisors on a curve are classically expressed by a divisor and its
residual divisor. What is established here is the identification
of these functors with specified multiplication Fitting schemes,
the canonical cokernel line \eqref{eq:hyperplane-canonical-line},
and the global-section transports. These are separate assertions
from the classical classification of geometric hyperplanes.

The faithful-action viewpoint is also classical; compare
\cite[Proposition 3.4]{Sistko} for subalgebras of matrix algebras.
Our conductor flags record the \emph{specified inclusion} into $B$,
the action on $B/E$, and its exact-rank scheme structure. A conductor
rank bound alone is not the flag equivalence, nor does either imply
unrestricted base change across an action-rank jump.
The one-variable codimension-two classification is already given
in \cite[Theorem 39]{GLTU}. We use it to locate the fibres of
Example~\ref{ex:conductor-rank-jump}, not as a new classification
of polynomial subalgebras. The quotient-tower statement concerns
arbitrary finite $B$ and its relative exact-rank functors.

In \cite[Example 4.16]{ABHS}, the two-generator jet calculation has
curvilinear target $k[\varepsilon]/\varepsilon^{m+1}$.
Our target in \eqref{eq:fat-algebra} has embedding dimension $e$.
The symmetric-power determinant identity in the proof is classical
linear algebra. Its role is to determine the full primary power of
the tangent determinant, uniformly in $e,h$, and then identify that
power with global section multiplication on $\Pj^e$ using
Theorem~\ref{thm:fat-conductor}. Merely knowing the polygenerator
open would not determine this exponent or the global degree range.
These comparisons identify the inputs and additional claims; they
are not a certificate of exhaustive priority against all literature.

\subsection*{Codimension, embedded incidences, and regularity}
The elementary dimension bound for an increasing sequence of subspaces
is not a scheme-theoretic stabilization theorem. The proof of
Theorem~\ref{thm:codimension-stabilization} retains the compressed
commutator relation over the coefficient ring before using polarized
Cayley--Hamilton. The polygenerator minors of \cite{ABHS} remain the
classical ambient construction. The assertions added here are the
codimension-sensitive degree bound, the full codimension-two matrix,
and its conductor incidence on nonreduced bases. We do not claim
an exhaustive priority comparison for every possible length bound.

The binary cubic for a rank-three algebra is the explicit square
condition for a line to lift to a subalgebra. No independent novelty
claim is made for that classical cubic-algebra calculation. Its role
in Proposition~\ref{prop:codim-two-incidence-fibres} is to describe
an exceptional fibre of the same incidence which transfers the
embedded local rings in Theorem~\ref{thm:embedded-multigenerator}.
The contact algebra $G_h$ and its Schur basis are classical inputs;
the mixed equations $bv,b^2u$, their displayed primary intersection,
and their realization as local rings of full multiplication failure
are proved separately.

Sidman's Theorems 1.3 and 1.8 in \cite{Sidman} supply the relationship
between regularity and saturation, and the ordinary-square estimate
$\operatorname{reg}(I^2)\le2\operatorname{reg}I$ for a finite
projective scheme. Sidman also records the earlier power bounds of
Chandler and of Geramita--Gimigliano--Pitteloud. None of these
regularity inequalities is claimed as new. Their use here is the
normal-layer identity \eqref{eq:normal-layer-fill} and the resulting
isomorphism of actual multiplication cokernels, including relative
families. This separation is particularly important: bounding the
regularity of an ideal is not by itself a theorem identifying the
Fitting scheme of a prescribed incomplete linear series.

Ballico's failure-locus work \cite{Ballico93,Ballico96} is a necessary
nearest-source comparison. The accessible statements of
\cite[Theorem 0.2, Proposition 2.2, Remark 2.3]{Ballico96} concern
failure on finite linear sections of a fixed completely embedded
variety. The complete theorem pages of \cite{Ballico93} have not
been obtained in this revision, so their exact relation to the
present scheme statements remains unverified. We make no inference
of non-anticipation from that access limitation. The results here
are submitted with complete proofs and explicit classical inputs;
priority of the resulting package is a separate question for source
comparison and editorial assessment.
